% =====================================================================
%  CHAPITRE 18 — LES 100 ERREURS TYPOIQUES DES FRANCOPHONES
% =====================================================================
\chapter{Les 100 erreurs qui trahissent un francophone}

Ce chapitre est un outil de diagnostic. Lisez-le une première fois en entier (30 minutes),
puis revenez-y une fois par mois. Chaque erreur est présentée avec sa forme fautive,
sa forme correcte et l'explication. \textbf{Ne travaillez qu'une erreur à la fois} :
choisissez-en trois par semaine, affichez-les, et corrigez-les dans tous vos écrits et
enregistrements jusqu'à l'automatisme.

\begin{tableau}[Erreurs 1 à 25 --- le verbe et ses temps]
\footnotesize
\begin{longtable}{@{}p{5.0cm}p{4.6cm}p{4.2cm}@{}}
\entete{\ko Forme fautive} & \entete{\ok Forme correcte} & \entete{Explication} \tabularnewline
\endfirsthead
\entete{\ko Faux} & \entete{\ok Juste} & \entete{Pourquoi} \tabularnewline
\endhead
\hline
\en{*I have 25 years.} & \en{I'm 25 (years old).} & L'âge se dit avec \en{be}. \tabularnewline
\en{*I have cold / hunger / fear.} & \en{I'm cold / hungry / afraid.} & Sensations avec \en{be} + adjectif. \tabularnewline
\en{*I am agree.} & \en{I agree.} & \en{agree} est un verbe, pas un adjectif. \tabularnewline
\en{*I am knowing him.} & \en{I know him.} & Verbes d'état jamais au \en{-ing}. \tabularnewline
\en{*I want that you come.} & \en{I want you to come.} & \en{want + personne + to}. \tabularnewline
\en{*He don't work.} & \en{He doesn't work.} & 3\textsuperscript{e} personne : \en{doesn't}. \tabularnewline
\en{*He doesn't has a car.} & \en{He doesn't have a car.} & Après \en{does}, base verbale. \tabularnewline
\en{*Did you went?} & \en{Did you go?} & Après \en{did}, base verbale. \tabularnewline
\en{*She don't like it.} & \en{She doesn't like it.} & Idem 3\textsuperscript{e} personne. \tabularnewline
\en{*I have seen him yesterday.} & \en{I saw him yesterday.} & Date précise \ra passé simple. \tabularnewline
\en{*I live here since 2015.} & \en{I've lived here since 2015.} & « Depuis » \ra \en{present perfect}. \tabularnewline
\en{*I am here since Monday.} & \en{I've been here since Monday.} & Idem avec \en{be}. \tabularnewline
\en{*I will call you when I will arrive.} & \en{I'll call you when I arrive.} & Présent après \en{when}. \tabularnewline
\en{*If I will have time\ldots} & \en{If I have time\ldots} & Jamais \en{will} après \en{if}. \tabularnewline
\en{*I used to go there when I was young} \emph{(correct)} \en{*I am used to go} & \en{I'm used to going.} & \en{be used to} + \en{-ing}. \tabularnewline
\en{*He can to swim.} & \en{He can swim.} & Pas de \en{to} après un modal. \tabularnewline
\en{*She cans / *he musts.} & \en{She can / he must.} & Les modaux ne prennent jamais \en{-s}. \tabularnewline
\en{*I mustn't come} \emph{(= interdit)} & \en{I don't have to come.} & « Pas obligé » \ra \en{don't have to}. \tabularnewline
\en{*I have gone to Paris twice.} & \en{I have been to Paris twice.} & Expérience \ra \en{been}. \tabularnewline
\en{*She's gone to the market} \emph{(elle est au marché)} & \en{She's gone to the market.} & Correct --- mais vérifiez le sens voulu. \tabularnewline
\en{*I didn't nothing say.} & \en{I didn't say anything.} & Une seule négation. \tabularnewline
\en{*I don't never go.} & \en{I never go.} & \en{never} suffit. \tabularnewline
\en{*I have make a mistake.} & \en{I have made a mistake.} & Participe passé, pas base. \tabularnewline
\en{*She has wrote it.} & \en{She has written it.} & 3\textsuperscript{e} forme. \tabularnewline
\en{*I didn't knew.} & \en{I didn't know.} & Base après \en{didn't}. \tabularnewline
\end{longtable}
\end{tableau}

\begin{tableau}[Erreurs 26 à 50 --- articles, noms, prépositions]
\footnotesize
\begin{longtable}{@{}p{5.0cm}p{4.6cm}p{4.2cm}@{}}
\entete{\ko Forme fautive} & \entete{\ok Forme correcte} & \entete{Explication} \tabularnewline
\endfirsthead
\entete{\ko Faux} & \entete{\ok Juste} & \entete{Pourquoi} \tabularnewline
\endhead
\hline
\en{*I like the music.} & \en{I like music.} & Généralité \ra pas d'article. \tabularnewline
\en{*The life is beautiful.} & \en{Life is beautiful.} & Concept abstrait général. \tabularnewline
\en{*I live in the Senegal.} & \en{I live in Senegal.} & Pays sans \en{the} (sauf pluriel/États). \tabularnewline
\en{*I go to the school} \emph{(étudier)} & \en{I go to school.} & Fonction \ra pas d'article. \tabularnewline
\en{*I had the breakfast at 7.} & \en{I had breakfast at 7.} & Repas sans article. \tabularnewline
\en{*She speaks the English.} & \en{She speaks English.} & Langues sans article. \tabularnewline
\en{*I play the football.} & \en{I play football.} & Sports sans article (mais \en{play \textbf{the} piano}). \tabularnewline
\en{*He is in the Monday.} & \en{He arrives on Monday.} & \en{on} + jour, sans article. \tabularnewline
\en{*I will see you in tomorrow.} & \en{I'll see you tomorrow.} & Jamais de préposition avant \en{tomorrow}. \tabularnewline
\en{*in the night} \emph{(en général)} & \en{at night} & \en{at night} mais \en{in the morning}. \tabularnewline
\en{*an university} & \en{a university} & \en{u} = \pron{iou} (consonne). \tabularnewline
\en{*a hour} & \en{an hour} & \en{h} muet. \tabularnewline
\en{*two informations} & \en{two pieces of information} & Indénombrable. \tabularnewline
\en{*many advices} & \en{a lot of advice} & Indénombrable. \tabularnewline
\en{*a news} & \en{a piece of news / some news} & \en{news} est singulier. \tabularnewline
\en{*the news are good} & \en{the news is good} & Accord au singulier. \tabularnewline
\en{*a hairs} & \en{hair} (chevelure) & \en{hair} est indénombrable ici. \tabularnewline
\en{*peoples} & \en{people} & Déjà pluriel. \tabularnewline
\en{*childrens} & \en{children} & Déjà pluriel. \tabularnewline
\en{*informations about} & \en{information about} & Indénombrable. \tabularnewline
\en{*depend of} & \en{depend on} & Préposition imposée. \tabularnewline
\en{*listen music} & \en{listen to music} & Préposition imposée. \tabularnewline
\en{*discuss about} & \en{discuss (something)} & Pas de préposition. \tabularnewline
\en{*married with} & \en{married to} & Préposition imposée. \tabularnewline
\en{*I have fear of dogs.} & \en{I'm afraid of dogs.} & Structure anglaise. \tabularnewline
\end{longtable}
\end{tableau}

\begin{tableau}[Erreurs 51 à 75 --- structure de phrase et vocabulaire]
\footnotesize
\begin{longtable}{@{}p{5.0cm}p{4.6cm}p{4.2cm}@{}}
\entete{\ko Forme fautive} & \entete{\ok Forme correcte} & \entete{Explication} \tabularnewline
\endfirsthead
\entete{\ko Faux} & \entete{\ok Juste} & \entete{Pourquoi} \tabularnewline
\endhead
\hline
\en{*Is cold today.} & \en{It's cold today.} & Sujet \en{it} obligatoire. \tabularnewline
\en{*Have a problem.} & \en{There's a problem.} & \en{il y a} = \en{there is}. \tabularnewline
\en{*It has two rooms.} & \en{It has got / There are two rooms.} & \en{il y a} $\neq$ \en{it has}. \tabularnewline
\en{*I am a Senegalese.} & \en{I'm Senegalese.} & Nationalité : adjectif sans article. \tabularnewline
\en{*She is a doctor woman.} & \en{She is a doctor.} & Pas de redondance. \tabularnewline
\en{*I said you something.} & \en{I told you something.} & \en{tell + personne}, \en{say + chose}. \tabularnewline
\en{*He explained me the problem.} & \en{He explained the problem to me.} & \en{explain} n'accepte pas le double complément direct. \tabularnewline
\en{*I demand you a question.} & \en{I'll ask you a question.} & \en{demand} = exiger. \tabularnewline
\en{*Make attention!} & \en{Be careful! / Watch out!} & Calque de « faire attention ». \tabularnewline
\en{*Make a photo.} & \en{Take a photo.} & Collocation imposée. \tabularnewline
\en{*Do a mistake.} & \en{Make a mistake.} & \en{make} pour créer, \en{do} pour exécuter. \tabularnewline
\en{*Do a party.} & \en{Have / throw a party.} & Collocation. \tabularnewline
\en{*Make sport.} & \en{Do sport / play sports.} & Collocation. \tabularnewline
\en{*I have a good relation with him.} & \en{I get on well with him.} & \en{relation} = relation amoureuse ou diplomatique. \tabularnewline
\en{*The people is nice here.} & \en{The people are nice here.} & \en{people} est pluriel. \tabularnewline
\en{*I am boring.} \emph{(je m'ennuie)} & \en{I'm bored.} & \en{-ed} = ressenti. \tabularnewline
\en{*The film was bored.} & \en{The film was boring.} & \en{-ing} = cause. \tabularnewline
\en{*I have 32 years old.} & \en{I'm 32.} & Redondance. \tabularnewline
\en{*It's difficult for me of understanding.} & \en{It's difficult for me to understand.} & \en{for + personne + to}. \tabularnewline
\en{*I think to go.} & \en{I'm thinking of going.} & \en{think of + -ing} (projeter). \tabularnewline
\en{*I look my phone.} & \en{I look at my phone.} & \en{look at} ; \en{watch} pour ce qui bouge. \tabularnewline
\en{*I see TV every night.} & \en{I watch TV every night.} & \en{watch} = attention prolongée. \tabularnewline
\en{*The weather is very hot} \emph{(correct, mais)} & \en{It's very hot.} & Plus naturel. \tabularnewline
\en{*How is it called?} & \en{What is it called?} & \en{what} pour un nom. \tabularnewline
\en{*How do you call this?} & \en{What do you call this?} & Même règle. \tabularnewline
\end{longtable}
\end{tableau}

\begin{tableau}[Erreurs 76 à 100 --- discours, prononciation, style]
\footnotesize
\begin{longtable}{@{}p{5.0cm}p{4.6cm}p{4.2cm}@{}}
\entete{\ko Forme fautive} & \entete{\ok Forme correcte} & \entete{Explication} \tabularnewline
\endfirsthead
\entete{\ko Faux} & \entete{\ok Juste} & \entete{Pourquoi} \tabularnewline
\endhead
\hline
\en{*He asked me where did I live.} & \en{He asked me where I lived.} & Question indirecte : ordre affirmatif. \tabularnewline
\en{*I don't know what does it mean.} & \en{I don't know what it means.} & Idem. \tabularnewline
\en{*The man who I met him.} & \en{The man (who) I met.} & Ne pas répéter le pronom. \tabularnewline
\en{*My brother, that lives in Dakar\ldots} & \en{My brother, who lives in Dakar\ldots} & \en{that} interdit en relative non définissante. \tabularnewline
\en{*I am living here since 2010.} & \en{I have been living here since 2010.} & Durée \ra \en{perfect}. \tabularnewline
\en{*She said me that\ldots} & \en{She told me that\ldots} & \en{say} ne prend pas de personne directement. \tabularnewline
\en{*I am since one hour here.} & \en{I have been here for an hour.} & \en{for} + durée. \tabularnewline
\en{*I search my keys.} & \en{I'm looking for my keys.} & \en{search} = fouiller un lieu. \tabularnewline
\en{*I will see you to tomorrow.} & \en{I'll see you tomorrow.} & Préposition inutile. \tabularnewline
\en{*Actually I'm working} \emph{(= actuellement)} & \en{Currently I'm working.} & \en{actually} = en fait. \tabularnewline
\en{*Eventually we'll meet} \emph{(= éventuellement)} & \en{We'll possibly meet.} & \en{eventually} = finalement. \tabularnewline
\en{*I visited my grandmother} \emph{(rendre visite)} & \en{I visited my grandmother.} \ok & Correct ici --- mais \en{visit a place} est plus fréquent ; pour les personnes, \en{see} est souvent plus naturel. \tabularnewline
\en{*I want a coffee} \emph{(au café)} & \en{I'd like a coffee, please.} & Registre : \en{want} est brutal. \tabularnewline
\en{*Give me the salt.} \emph{(à table, adulte)} & \en{Could you pass me the salt, please?} & Politesse obligatoire entre adultes. \tabularnewline
\en{*I have the possibility to come.} & \en{I can come.} / \en{I have the chance to come.} & \en{possibility} = faisabilité technique. \tabularnewline
\en{*It is necessary that you come.} & \en{You need to come.} / \en{You must come.} & Éviter les structures lourdes. \tabularnewline
\en{*I suggest you to come.} & \en{I suggest coming.} / \en{I suggest that you come.} & \en{suggest} refuse \en{suggest sb to do}. \tabularnewline
\en{*He explained me everything} & \en{He explained everything to me.} & Structure imposée. \tabularnewline
\en{*I prefer coffee than tea.} & \en{I prefer coffee to tea.} & \en{prefer A to B}. \tabularnewline
\en{*I am agree with you} & \en{I agree with you.} & Déjà vu, mais faute tenace. \tabularnewline
\en{*This is more better.} & \en{This is better.} & Jamais \en{more} + comparatif. \tabularnewline
\en{*He runs more fast than me.} & \en{He runs faster than me.} & Comparatif en \en{-er}. \tabularnewline
\en{*I have so much friends.} & \en{I have so many friends.} & \en{many} + dénombrable. \tabularnewline
\en{*There is many people.} & \en{There are many people.} & Accord avec le pluriel. \tabularnewline
\en{*I look forward to hear from you.} & \en{I look forward to hearing from you.} & \en{to} est une préposition ici. \tabularnewline
\end{longtable}
\end{tableau}

\section{Les cinq fautes de prononciation à corriger en premier}
\begin{enumerate}
  \item \textbf{Le \en{h} avalé} : \en{he, him, house, happy} --- soufflez.
  \item \textbf{Le \en{th} transformé} : \en{think} $\neq$ \en{sink} ; \en{this} $\neq$ \en{dis}.
  \item \textbf{Le \en{w} dit \en{v}} : \en{we} $\neq$ \en{ve} ; \en{water} $\neq$ \en{vater}.
  \item \textbf{Le \en{e} final ajouté} : \en{book} \pron{bouk}, pas \pron{bou-keu}.
  \item \textbf{L'accent tonique plat} : \en{baNAna}, \en{inforMAtion}, \en{toDAY}.
\end{enumerate}

\section{Auto-diagnostic : 40 phrases à corriger}
Couvrez la colonne de droite, corrigez les 40 phrases, puis comparez. Score sur 40 :
\textbf{35 et plus} = niveau B1 solide sur les structures de base ; \textbf{25-34} =
niveau A2, travaillez trois erreurs par semaine ; \textbf{moins de 25} = reprenez les
chapitres 8 et 9 en priorité.

\begin{tableau}[Test de diagnostic]
\footnotesize
\begin{longtable}{@{}p{0.9cm}p{6.2cm}p{6.7cm}@{}}
\entete{\#} & \entete{Phrase à corriger} & \entete{Correction} \tabularnewline
\endfirsthead
\entete{\#} & \entete{Phrase à corriger} & \entete{Correction} \tabularnewline
\endhead
1 & \en{I have 40 years.} & \en{I'm 40 (years old).} \tabularnewline
2 & \en{She don't like fish.} & \en{She doesn't like fish.} \tabularnewline
3 & \en{I am agree with you.} & \en{I agree with you.} \tabularnewline
4 & \en{He can to help us.} & \en{He can help us.} \tabularnewline
5 & \en{I live here since five years.} & \en{I've lived here for five years.} \tabularnewline
6 & \en{Yesterday I have seen Ali.} & \en{Yesterday I saw Ali.} \tabularnewline
7 & \en{I will call you when I will arrive.} & \en{I'll call you when I arrive.} \tabularnewline
8 & \en{If it will rain, we'll stay.} & \en{If it rains, we'll stay.} \tabularnewline
9 & \en{I have cold.} & \en{I'm cold.} \tabularnewline
10 & \en{The life is hard.} & \en{Life is hard.} \tabularnewline
11 & \en{I like the football.} & \en{I like football.} \tabularnewline
12 & \en{He explained me the rule.} & \en{He explained the rule to me.} \tabularnewline
13 & \en{She said me she was tired.} & \en{She told me she was tired.} \tabularnewline
14 & \en{I am boring today.} & \en{I'm bored today.} \tabularnewline
15 & \en{Make attention to the step.} & \en{Watch your step.} / \en{Be careful.} \tabularnewline
16 & \en{I did a mistake.} & \en{I made a mistake.} \tabularnewline
17 & \en{He asked me where do I live.} & \en{He asked me where I live(d).} \tabularnewline
18 & \en{The man who I met him yesterday.} & \en{The man (who) I met yesterday.} \tabularnewline
19 & \en{My sister, that works in a bank\ldots} & \en{My sister, who works in a bank\ldots} \tabularnewline
20 & \en{I search my glasses.} & \en{I'm looking for my glasses.} \tabularnewline
21 & \en{She is married with a teacher.} & \en{She is married to a teacher.} \tabularnewline
22 & \en{I look TV every evening.} & \en{I watch TV every evening.} \tabularnewline
23 & \en{Actually I'm busy.} \emph{(= actuellement)} & \en{Currently / At the moment I'm busy.} \tabularnewline
24 & \en{I prefer tea than coffee.} & \en{I prefer tea to coffee.} \tabularnewline
25 & \en{This is more easy.} & \en{This is easier.} \tabularnewline
26 & \en{I have so much work.} \emph{(correct !)} & \en{I have so much work.} --- \en{work} est indénombrable, donc \en{much} est juste. \tabularnewline
27 & \en{There is many problems.} & \en{There are many problems.} \tabularnewline
28 & \en{I look forward to meet you.} & \en{I look forward to meeting you.} \tabularnewline
29 & \en{I want that you help me.} & \en{I want you to help me.} \tabularnewline
30 & \en{I am knowing the answer.} & \en{I know the answer.} \tabularnewline
31 & \en{I didn't went there.} & \en{I didn't go there.} \tabularnewline
32 & \en{She has wrote a letter.} & \en{She has written a letter.} \tabularnewline
33 & \en{You mustn't pay} \emph{(= pas obligé)} & \en{You don't have to pay.} \tabularnewline
34 & \en{I used to go there} \emph{(= j'ai l'habitude d'y aller)} & \en{I usually go there.} \tabularnewline
35 & \en{I'm used to wake up early.} & \en{I'm used to waking up early.} \tabularnewline
36 & \en{I have been in Paris two times.} & \en{I have been to Paris twice.} \tabularnewline
37 & \en{The news are bad.} & \en{The news is bad.} \tabularnewline
38 & \en{He gave me an advice.} & \en{He gave me a piece of advice.} \tabularnewline
39 & \en{I'll see you in tomorrow.} & \en{I'll see you tomorrow.} \tabularnewline
40 & \en{How is it called in English?} & \en{What is it called in English?} \tabularnewline
\end{longtable}
\end{tableau}

\begin{mission}[Chapitre 18 --- trois erreurs par semaine]
Chaque lundi, choisissez \textbf{trois} erreurs de ce chapitre (pas plus). Écrivez-les sur
un papier que vous gardez dans votre poche ou sur votre bureau. Pendant toute la semaine,
chaque fois que vous parlez ou écrivez en anglais, vérifiez ces trois points. Le dimanche,
faites un test : écrivez 10 phrases en utilisant les trois structures surveillées.
En 33 semaines, vous aurez éliminé les 100 erreurs --- et personne ne devinera plus que
vous êtes francophone en vous lisant.
\end{mission}
